
\subsubsection{6.1 Key Findings}

\textbf{Two mechanistically distinct pathways to WGA improvement.} DFR achieves WGA = 0.91 by recalibrating the head while leaving backbone features --- with their full spurious background encoding (ERM probe accuracy $\approx$ 0.984) --- entirely unchanged. GroupDRO achieves WGA = 0.88 by modifying backbone representations such that background information is significantly less linearly extractable (p = 0.0039, d = 6.48). These are two distinct robustification pathways, not two points on a continuum.

A non-obvious implication is that the backbone's spurious encoding level is neither necessary nor sufficient for WGA improvement in isolation. DFR demonstrates that high spurious encoding in the backbone does not prevent high WGA if the head is correctly calibrated. GroupDRO demonstrates that reducing backbone spurious encoding also yields high WGA. The mechanisms are alternative solutions to the same problem.

\textbf{GroupDRO backbone modification is substantial and consistent.} The backbone/head L2 ratio of 6.47--7.03 indicates that GroupDRO's group-reweighted gradient signal primarily reshapes backbone weights. The DFR negative control (backbone/head ratio = 0.000 for all seeds) eliminates checkpoint-specific artifact confounds: any backbone modification observed for GroupDRO is attributable to training dynamics, not checkpoint artifacts.

\textbf{Why is the effect size so large (d = 6.48)?} The large Cohen's d reflects high stability of probe accuracy estimates rather than an implausibly large biological effect. When probe accuracy is computed on N=5,794 test samples, the estimate is highly precise across seeds. The per-pair SD of $ERM-GroupDRO$ differences is approximately 0.00476, producing a large standardized effect for a consistent 0.03-unit shift that is stable across all three seed pairs. The DFR negative control (cosine similarity = 1.000000) confirms the probe is detecting genuine backbone differences rather than methodological artifact.

\textbf{Connection to the literature.} GroupDRO implicitly reduces background linear decodability, aligning with Park et al. [2025] (SCER explicitly regularizes spurious feature directions). H-M2 results show the backbone changes substantially (ratio 6.47--7.03), partially contesting the view of Izmailov et al. [2022] that GroupDRO's advantage is ''primarily in the head'' --- though backbone modification does not preclude head-level contributions to WGA. H-M2 weight-difference analysis corroborates Raymond et al. [2026, preprint]; this finding stands on its own pre-registered evidence independent of that citation.

\subsubsection{6.2 Limitations}

\textbf{L1: n=3 seeds.} One-sided t-test with n=3 has approximately 55\% power for d=0.8. The primary finding (d=6.48) requires only n=2 for 80\% power at $\alpha =0.05$ one-sided and is unaffected by the small-n limitation. H-M2's linear probe result (p=0.0526, d=1.64) is borderline SUGGESTIVE; additional seeds would resolve this.

\textbf{L2: Single dataset and architecture.} All experiments use Waterbirds WILDS with ResNet-50. Generalization to CelebA, MultiNLI, ViT, or DINO architectures is an open question not addressed by this study.

\textbf{L3: H-P2 SUGGESTIVE at n=9.} The pre-registered CONFIRMED threshold (CI upper bound < 0) is not met. Root cause: method-level WGA constants from Izmailov et al. [2022] create within-method collinearity that inflates bootstrap CI width. The ERM+GroupDRO ablation (r=$-0.755$, p=0.041) is CONFIRMED. Obtaining per-seed WGA values would likely resolve this limitation.

\textbf{L4: Decodability $\neq$ mechanism.} Reduced probe accuracy establishes that background information is less linearly extractable from GroupDRO features; it does not establish \textit{why}. Two competing explanations remain: (a) spurious feature suppression (GroupDRO actively reduces the strength of background-predictive directions), or (b) feature dilution (increased feature diversity reduces the proportion of spurious information). Both produce identical probe accuracy signatures. Distinguishing them requires directional subspace analysis (e.g., SCER-style spurious subspace probing \citep{park2025scer}).

\textbf{L5: Core attribute probe not measured.} Bird species (core attribute) probe accuracy from GroupDRO features was not measured in H-M3. Whether GroupDRO preserves core-feature encoding is an important open validation question.

Establishing causality would require intervention studies (e.g., ablating minority upweighting while holding architecture fixed). The present results establish consistent mechanistic co-occurrence, not causal necessity.

\subsubsection{6.3 Practical Implications}

The diagnostic framework --- spurious attribute linear probe accuracy on frozen backbone features --- is low-cost (no new training, forward-pass only), reproducible (standard sklearn tools, public checkpoints), and directly interpretable. It produces a backbone spurious encoding ''fingerprint'' for any robustification method with publicly available checkpoints.

Understanding the backbone-vs-head distinction informs method selection: when group labels for head retraining are available, DFR's simpler head-only intervention achieves the highest WGA in the Waterbirds setting. When backbone-level change is desired for transfer learning, or when group annotations are unavailable at inference time, GroupDRO's backbone modification offers a substantively different form of robustification.

